%==============================================================================
\chapter{السياق وتعريف المشكلة}
%==============================================================================

\section{عرض عام}
يهدف مشروع \textenglish{Zümm} إلى تصميم نظام معلومات مخصّص لإدارة ومتابعة
أنشطة تربية النحل. تعتمد تربية النحل الحديثة على المراقبة المنتظمة للعديد من
الخلايا، الموزّعة غالبًا على عدة مواقع جغرافية والمُستغَلّة من قِبَل مربّين
مختلفين. تمثّل متابعة صحة الطوائف وإنتاجيتها وعمليات الصيانة عبئًا تنظيميًا
كبيرًا، يُدار اليوم بطريقة يدوية ومشتّتة إلى حدّ بعيد.

\section{الإشكالية}
في غياب أداة مركزية، يواجه مستغلّو تربية النحل عدة صعوبات:
\begin{itemize}
  \item غياب التتبّع المنظّم لزيارات الخلايا وعمليات التفتيش
  \item صعوبة تخطيط وتنسيق عمل أعوان تربية النحل
  \item نقص الرؤية حول الإنتاج (العسل، الوزن) وحول مردودية كل موقع أو خلية
  \item الكشف المتأخّر عن الشذوذات (تراجع العدد، انخفاض الإنتاج، فتح غير مبرمَج لخلية)
  \item استحالة الاستغلال الآلي للقياسات الواردة من مستشعرات غير متجانسة.
\end{itemize}

\section{الغاية}
يجب أن يقدّم النظام استجابة منظّمة في شقّين متكاملين:
\begin{enumerate}
  \item \textbf{إدارة يدوية} لعمليات المناولة وتخطيطها ومتابعتها (الجزء الأول،
        موضوع هذا المشروع)
  \item \textbf{إشراف آلي عن بُعد} على مؤشّرات الأداء، تغذّيه المستشعرات، مُعَدّ
        ليُطوَّر في إطار مشروع لاحق، لكن على النظام أن يستبق دمجه.
\end{enumerate}
